% !TEX root = chapter-standalone.tex

\section[Analysis]{Analysis}
\label{sec:cats-of-analysis}

In analysis one studies functions between the Euclidean spaces~$\reals^n$ that satisfy some regularity property: being continuous, being differentiable, not growing too fast, and so on.
Each of these properties gives rise to a category via the following recipe: the objects are the spaces~$\reals^n$, and the morphisms are the functions with the given property.
For this recipe to work, two facts must hold for the property in question: identity functions must have it, and it must be preserved under composition of functions.
We spell this out for three of the most common properties.

\subsection{Continuous functions}

\begin{definition}[Continuous function]
    \label{def:continuous-function}
    A function~$\mapa \colon \reals^n \sto \reals^m$ is \emph{continuous} at a point~$\elc \setin \reals^n$ if
    \begin{equation}
        \lim_{\ela \to \elc} \mapa(\ela) = \mapa(\elc),
    \end{equation}
    and~$\mapa$ is \emph{continuous} if it is continuous at every point of~$\reals^n$.
\end{definition}

Identity functions are continuous, and it is a basic fact of analysis that the composition of two continuous functions is again continuous.
Hence we may define:

\begin{ctdefinition}[Category of Euclidean spaces and continuous functions]
    \SYNDEF{category of Euclidean spaces and continuous functions}
    \label{def:cat-Cont}
    The category~\Cont is:
    \begin{enumerate}
        \item \emph{Objects}: the Euclidean spaces~$\reals^n$, for~$n \setin \natnumbers$.
        \item \emph{Morphisms}: $\HomSet{\Cont}{\reals^n}{\reals^m}$ is the set of continuous functions~$\reals^n \sto \reals^m$.
        \item \emph{Composition}: the usual composition of functions.
        \item \emph{Identity morphisms}: identity functions.
    \end{enumerate}
\end{ctdefinition}

\subsection{Smooth functions}

\begin{definition}[Smooth function]
    \label{def:smooth-function}
    A function~$\mapa \colon \reals^n \sto \reals^m$ is \emph{smooth} (or \emph{infinitely differentiable}) if all of its partial derivatives, of all orders, exist and are continuous.
\end{definition}

Identity functions are smooth, and by the chain rule the composition of two smooth functions is smooth.
Hence we may define:

\begin{ctdefinition}[Category of Euclidean spaces and smooth functions]
    \SYNDEF{category of Euclidean spaces and smooth functions}
    \label{def:cat-Smooth}
    The category~\Smooth is:
    \begin{enumerate}
        \item \emph{Objects}: the Euclidean spaces~$\reals^n$, for~$n \setin \natnumbers$.
        \item \emph{Morphisms}: $\HomSet{\Smooth}{\reals^n}{\reals^m}$ is the set of smooth functions~$\reals^n \sto \reals^m$.
        \item \emph{Composition}: the usual composition of functions.
        \item \emph{Identity morphisms}: identity functions.
    \end{enumerate}
\end{ctdefinition}

\subsection{Lipschitz functions}

\begin{definition}[Lipschitz function]
    \label{def:Lipschitz}
    A function~$\mapa \colon \reals^n \sto \reals^m$ is \emph{Lipschitz} (or \emph{Lipschitz continuous}) if there exists a constant~$\kappa \geq 0$ such that
    \begin{equation}
        \lVert \mapa(\ela) - \mapa(\elb) \rVert \leq \kappa \, \lVert \ela - \elb \rVert
        \quad \text{for all } \ela, \elb \setin \reals^n,
    \end{equation}
    where~$\lVert \,\cdot\, \rVert$ denotes the Euclidean norm.
\end{definition}

Identity functions are Lipschitz (with~$\kappa = 1$).
If~$\mapa$ is Lipschitz with constant~$\kappa_\mapa$ and~$\mapb$ is Lipschitz with constant~$\kappa_\mapb$, then~$\mapa \mthen \mapb$ is Lipschitz with constant~$\kappa_\mapa \kappa_\mapb$:
\begin{equation}
    \lVert \mapb(\mapa(\ela)) - \mapb(\mapa(\elb)) \rVert
    \leq \kappa_\mapb \, \lVert \mapa(\ela) - \mapa(\elb) \rVert
    \leq \kappa_\mapb \kappa_\mapa \, \lVert \ela - \elb \rVert.
\end{equation}
Hence we may define:

\begin{ctdefinition}[Category of Euclidean spaces and Lipschitz functions]
    \SYNDEF{category of Euclidean spaces and Lipschitz functions}
    \label{def:cat-Lip}
    The category~\Lip is:
    \begin{enumerate}
        \item \emph{Objects}: the Euclidean spaces~$\reals^n$, for~$n \setin \natnumbers$.
        \item \emph{Morphisms}: $\HomSet{\Lip}{\reals^n}{\reals^m}$ is the set of Lipschitz functions~$\reals^n \sto \reals^m$.
        \item \emph{Composition}: the usual composition of functions.
        \item \emph{Identity morphisms}: identity functions.
    \end{enumerate}
\end{ctdefinition}

\begin{remark}
    Smooth functions and Lipschitz functions are in particular continuous, so every morphism of~\Smooth or of~\Lip is also a morphism of~\Cont.
    We will say in \cref{sec:subcategories} that \Smooth and~\Lip are \emph{subcategories} of~\Cont, and that all three are subcategories of~\Set.
\end{remark}

\begin{remark}
    The three definitions above use progressively less of the structure of~$\reals^n$: continuity only needs a notion of limit (a \emph{topology}), the Lipschitz condition only needs a notion of distance (a \emph{metric}), and smoothness needs derivatives.
    Accordingly, each of the three categories is a small part of a larger one studied in mathematics, namely the categories of topological spaces and continuous maps, of metric spaces and Lipschitz maps, and of smooth manifolds and smooth maps.
    We will not need these larger categories in this book.
\end{remark}

\subsection{Exercises}

\begin{gradedexercise}[\exname{PointedEuclideanSpaces}]
    \label{ex:pointed-euclidean-spaces}

    In this exercise we will define a category $\pEuc$ of ``pointed Euclidean spaces'' and your task is to check that it is in fact a category.

    The objects of $\pEuc$ are pairs $\tup{\reals^n, \ela}$ where $\ela \setin \reals^n$ (and $n \setin \natnumbers$ varies).

    A morphism $\mora \colon \tup{\reals^n, \ela} \mto_{\pEuc} \tup{\reals^m, \elb}$ is a smooth (infinitely differentiable) function $\mora \colon \reals^n \mto \reals^m$ such that $\mora(\ela) = \elb$.

    Composition is given by composition of smooth functions, and identity morphisms are identity functions.

\end{gradedexercise}

\solutionof{PointedEuclideanSpaces}
